% Full bibliography replaces the old proposal-only references and history.
% Explicit labels preserve the Introduction's original [1]--[9] numbering.
\section*{References}
\setsource{}
\begin{frame}[c,label=references-nuclear]{References: ML in Nuclear Physics}
\footnotesize
\begin{thebibliography}{99}
\setlength{\itemsep}{1.1em}
\bibitem[1]{lasseri2020}
R. Lasseri, D. Regnier, J. Ebran, and A. Penon.
\href{https://doi.org/10.1103/PhysRevLett.124.162502}{Taming Nuclear Complexity
with a Committee of Multilayer Neural Networks.}
\textit{Phys. Rev. Lett.} \textbf{124}, 162502 (2020).
\bibitem[2]{barsotti2020}
R. Barsotti and M. R. Shepherd.
\href{https://doi.org/10.1088/1748-0221/15/05/P05021}{Using machine learning to
separate hadronic and electromagnetic interactions in the GlueX forward calorimeter.}
\textit{JINST} \textbf{15}, P05021 (2020).
\bibitem[3]{fanelli2020}
C. Fanelli and J. Pomponi.
\href{https://doi.org/10.1088/2632-2153/ab845a}{DeepRICH: Learning deeply
Cherenkov detectors.}
\textit{Mach. Learn.: Sci. Technol.} \textbf{1}, 015010 (2020).
\bibitem[4]{scheinker2019}
A. Scheinker et al.
\href{https://doi.org/10.1103/PhysRevAccelBeams.22.082802}{Model-independent
tuning for maximizing free electron laser pulse energy.}
\textit{Phys. Rev. Accel. Beams} \textbf{22}, 082802 (2019).
\end{thebibliography}
\end{frame}

\setsource{}
\begin{frame}[c,label=references-astro]{References: ML in Astrophysics}
\footnotesize
\begin{thebibliography}{99}
\setlength{\itemsep}{0.85em}
\bibitem[5]{moriwaki2023}
K. Moriwaki, T. Nishimichi, and N. Yoshida.
Machine learning for observational cosmology.
\textit{Rep. Prog. Phys.} \textbf{86}, 076901 (2023).
\bibitem[6]{fraser2023}
T. S. Fraser, R. Tojeiro, and H. G. Chittenden.
Applying unsupervised learning to resolve evolutionary histories and explore
the galaxy-halo connection in IllustrisTNG.
\textit{MNRAS} \textbf{522}, 5758--5774 (2023).
\bibitem[7]{arimura2023}
R. Arimura.
\href{https://doi.org/10.1051/epjconf/202328302011}{A machine learning approach
for mass composition analysis with TALE-SD data.}
\textit{EPJ Web Conf.} \textbf{283}, 02011 (2023).
\bibitem[8]{denhartogh2023}
J. W. den Hartogh et al.
Barium stars as tracers of s-process nucleosynthesis in AGB stars.
II. Using machine learning techniques on 169 stars.
\textit{Astron. Astrophys.} \textbf{672}, A143 (2023).
\bibitem[9]{smith2024}
M. S. Smith and D. Lu.
\href{https://doi.org/10.3389/fspas.2024.1494439}{Machine learning opportunities
for nucleosynthesis studies.}
\textit{Front. Astron. Space Sci.} \textbf{11}, 1494439 (2024).
\end{thebibliography}
\end{frame}

\setsource{}
\begin{frame}[c,label=references-agents]{References: Scientific Reasoning and Agents}
\footnotesize
\begin{thebibliography}{99}
\setlength{\itemsep}{1.1em}
\bibitem[10]{robin2026}
A. E. Ghareeb et al.
\href{https://doi.org/10.1038/s41586-026-10652-y}{A multi-agent system for
automating scientific discovery (Robin).}
\textit{Nature} (2026). DOI: 10.1038/s41586-026-10652-y.
\bibitem[11]{coscientist2025}
J. Gottweis et al.
\href{https://arxiv.org/abs/2502.18864}{Towards an AI co-scientist.}
arXiv:2502.18864 (2025).
Journal version: \href{https://doi.org/10.1038/s41586-026-10644-y}{Accelerating
scientific discovery with Co-Scientist}, \textit{Nature} (2026).
\bibitem[12]{xunzi2026}
X. Huang et al.
\href{https://doi.org/10.1038/s41551-026-01769-6}{XunZi, an AI biologist,
reveals disease-modifying targets.}
\textit{Nature Biomedical Engineering} (2026).
DOI: 10.1038/s41551-026-01769-6.
\bibitem[13]{ideascientist2026}
J. Liu et al.
\href{https://arxiv.org/abs/2610.04074}{IdeaScientist: Orchestrating Agents
for Grounded Scientific Ideation.}
arXiv:2610.04074 (2026).
\end{thebibliography}
\end{frame}

\setsource{}
\begin{frame}[c,label=references-workflows]{References: Code and Experimental Workflows}
\footnotesize
\begin{thebibliography}{99}
\setlength{\itemsep}{1.1em}
\bibitem[14]{sakana2026}
C. Lu et al.
\href{https://doi.org/10.1038/s41586-026-10265-5}{Towards end-to-end automation
of AI research.}
\textit{Nature} \textbf{651}, 914--919 (2026).
\bibitem[15]{paper2agent2026}
J. Miao, J. R. Davis, Y. Zhang, J. K. Pritchard, and J. Zou.
\href{https://doi.org/10.1038/s41586-026-11044-y}{Reimagining research papers
as interactive and reliable AI agents.}
\textit{Nature} (2026). DOI: 10.1038/s41586-026-11044-y.
\bibitem[16]{era2026}
E. Ayg\"un et al.
\href{https://doi.org/10.1038/s41586-026-10658-6}{An AI system to help
scientists write expert-level empirical software.}
\textit{Nature} (2026). DOI: 10.1038/s41586-026-10658-6.
\bibitem[17]{xray2026}
Z. Chen et al.
\href{https://doi.org/10.1038/s42256-026-01261-5}{An agentic artificially
intelligent X-ray scientist.}
\textit{Nature Machine Intelligence} \textbf{8}, 1075--1086 (2026).
\end{thebibliography}
\end{frame}

\setsource{}
\begin{frame}[c,label=references-experiments]{References: Nuclear Experiments and Deep Field}
\footnotesize
\begin{thebibliography}{99}
\setlength{\itemsep}{1.5em}
\bibitem[18]{miskovich2022}
S. A. Miskovich et al.
\href{https://doi.org/10.1103/PhysRevAccelBeams.25.044601}{Online Bayesian
optimization for a recoil mass separator.}
\textit{Phys. Rev. Accel. Beams} \textbf{25}, 044601 (2022).
\bibitem[19]{tsintari2025}
P. Tsintari et al.
\href{https://doi.org/10.1103/PhysRevResearch.7.013074}{Machine learning
enabled measurements of astrophysical $(p,n)$ reactions with the SECAR
recoil separator.}
\textit{Phys. Rev. Research} \textbf{7}, 013074 (2025).
\bibitem[20]{hubble}
NASA Science.
\href{https://science.nasa.gov/mission/hubble/science/universe-uncovered/hubble-deep-fields/}{Hubble's Deep Fields.}
Hubble Space Telescope science resource; source of the Deep Field image
and observing history.
\end{thebibliography}
\end{frame}


